后量子安全消息协议的分析通常聚焦单次握手、单条消息或单个会话，而异步接收接口还会把多个发送方关联输入组合为一次处理。单个条目的来源或安全性质只能刻画该条目本身，不能自动确定同批多个条目之间应保持的关系。安全消息的会话组合研究也表明，组合边界可能承担独立于组件性质的安全义务\cite{cremers2023session}。因此，批次准入中的条目关系需要单独建模。

\kwaay 是一种后量子、类似X3DH的可否认认证密钥交换协议\cite{collins2024kwaay}。其\batchreceive 操作在共享接收方状态下处理多个发送方关联输入；原文第5.1节（第21页）明确要求同一次调用中的元素对应不同参与方\cite{collins2024kwaayfull}。本文将该要求称为批内参与方互异条件（distinct-party condition）。研究对象不是协议是否遗漏这一条件，而是该条件在批次内部维持何种关系，以及条目各自具有合法来源时，目标关系是否能够由其他限制推出。

Signal、PQXDH等安全消息协议的形式化分析主要考察握手、会话和密钥安全性质\cite{cohngordon2020signal,bhargavan2024pqxdh}，Tamarin则为安全协议的符号执行轨迹与性质验证提供了通用方法\cite{meier2013tamarin}。身份名称、单射一致性和会话组合方面的研究进一步说明，参与方表示、事件对应关系与组合层性质应分别界定\cite{lupetti2006names,lowe1997hierarchy,cremers2023session}。这些工作提供了分析工具和概念基础，但没有直接回答\batchreceive 中消息、发送实例与参与方三个比较维度之间的关系。

若同一参与方能够通过不同发送实例产生不同消息，那么批内消息互异性、发送与接受事件的单射关系以及批内参与方互异性便具有不同的比较对象。本文据此考察：移除目标条件后，重复参与方条目能否进入同一批次并产生接受事件；一个与完整元组匹配的发送来源能否对应同批两个接受事件；消息层面的互异约束能否排除同一参与方的不同消息共同进入批次。研究边界限定在固定两槽的批次准入抽象，不把模型结果扩大为完整K-Waay攻击或部署缺陷。

为控制流程差异，本文在相同的发送、网络暴露、两槽收集、来源匹配和顺序处理生命周期上，仅改变准入约束。无互异约束模型用于观察移除目标条件后的行为；消息互异约束模型检验消息关系能否替代参与方关系；参与方互异约束模型作为直接施加目标约束的对照模型（positive control），下文简称目标对照。验证结果显示：第一种模型允许一个匹配发送来源对应同批两个接受事件；第二种模型虽满足批内消息互异性和批内来源单射性，同一参与方产生的不同消息仍可在同一批次中被接受；第三种模型保持批内参与方互异性，且有效的不同参与方批次可达。

本文的主要贡献如下：
\begin{enumerate}[label=\arabic*)]
  \item 形式化刻画\batchreceive 的批内参与方互异条件，明确区分参与方、发送实例和消息三个比较维度；
  \item 在共同两槽生命周期上构造三种受控批次准入模型，比较消息互异约束与参与方互异约束的作用，并以发送来源对应性和批内来源单射性刻画事件级差异；
  \item 给出机器验证支持的非替代性结论：批内消息互异性与批内来源单射性即使同时成立，也不足以推出批内参与方互异性；并通过目标对照确认批内参与方互异性成立，且有效的不同参与方批次仍然可达。
\end{enumerate}

